\begin{frame}{FAQ (1/2)}
  \begin{itemize}
    \item \textbf{Q. 마르코프 성질이 실제 문제에서 항상
      성립하나요?}
      \begin{itemize}
        \item 엄밀하게는 \textbf{거의 항상 근사}.
        \item 예: 자율주행에서 ``카메라 한 장''만 상태로 쓰면
          \textbf{속도·가속도 정보가 빠져} 마르코프 성질이 깨진다
          (다음 상태를 예측하려면 ``지금 얼마나 빠르게
          움직이고 있었는가''가 필요한데, 정지 이미지 한 장에는
          그 정보가 없다).
        \item 실전: \textbf{최근 몇 프레임을 함께 묶어}
          ``상태''로 재정의하거나, \textbf{신경망이 과거 정보를
          요약한 은닉 상태}를 상태의 일부로 쓰는 방식
          (Week 11 RNN 은닉 상태와 같은 아이디어)으로 다시 성립.
      \end{itemize}
    \item \textbf{Q. 밴딧(Week 2)도 MDP의 특수한 경우인가요?}
      \begin{itemize}
        \item \textbf{그렇다} --- $\mathcal{S} = \{s_0\}$ (상태
          하나)에서 어떤 행동을 해도 \textbf{항상 같은 상태로
          ``전이''}하는 특수한 MDP.
        \item Week 2.3에서 밴딧과 MDP를 ``별도''로 소개했지만,
          이렇게 보면 이번 장의 \textbf{모든 정의}(가치함수,
          벨만방정식)가 밴딧에도 \textbf{퇴화된 형태로} 그대로
          적용.
      \end{itemize}
  \end{itemize}
\end{frame}
